\chapter{Regresión lineal}

\section{Apuntes del Tema 3: Regresión Lineal}

\subsection{1. Introducción y Planteamiento del Modelo}
\begin{itemize}
	\item \textbf{Definición:} Modelo paramétrico supervisado para predecir un valor objetivo continuo $y \in \mathbb{R}$ a partir de un vector de características $x \in \mathbb{R}^d$.
	\item \textbf{Ecuación del Modelo:} $\hat{y} = x^\top w + b$.
\end{itemize}

\subsection{2. Función de Pérdida (MSE)}
\begin{itemize}
	\item \textbf{MSE:} $\mathcal{L}(w, b) = \frac{1}{n} \Vert y - (Xw + b\mathbf{1}) \Vert_2^2$.
	\item \textbf{Derivación Probabilística:} Surge al asumir ruido gaussiano $\epsilon \sim \mathcal{N}(0, \sigma^2)$, donde maximizar la verosimilitud equivale a minimizar el MSE.
\end{itemize}

\subsection{3. Optimización}
\begin{itemize}
	\item \textbf{Analítica:} $\tilde{w} = (\tilde{X}^\top \tilde{X})^{-1} \tilde{X}^\top y$. (Costoso $O(d^3)$).
	\item \textbf{Descenso de Gradiente:} Actualización iterativa basada en el gradiente del MSE.
\end{itemize}

\subsection{4. Escalado y Expansión Polinomial}
\begin{itemize}
	\item \textbf{Estandarización:} Fundamental para evitar el "zig-zag" en el descenso de gradiente.
	\item \textbf{Expansión Polinomial:} Permite capturar fronteras no lineales conservando la linealidad en parámetros, aumentando el riesgo de \emph{overfitting}.
\end{itemize}

\section{Conexión con la práctica}
\subsection{Fuente: \texttt{03\_01\_Linear\_regression.ipynb}}
\begin{lstlisting}[language=Python]
	# Implementacion basica con PyTorch
	x = np.random.uniform(-3, 3, 200)
	y = w_true*x + b_true + noise
\end{lstlisting}